% sec:geo-photon-rings — Photon Rings and Bound Orbits
\section{Photon Rings and Bound Orbits\starmark{}}
\label{sec:geo-photon-rings}

Every image of a black hole --- from a bare Schwarzschild shadow to a
Kerr disc rendering with full radiative transfer --- contains a
characteristic bright ring near the shadow boundary.  This feature is
not a property of the accretion disc or the emission model: it is a
direct imprint of the unstable circular photon orbits that define the
photon sphere (\cref{sec:schw-orbits}).  Photons with impact parameter
$b$ close to the critical value $b_c$ wind multiple times around the
black hole before escaping, sweeping up light from every direction and
creating a feature whose structure encodes the spacetime geometry
independently of the astrophysical environment.  What makes the photon
ring remarkable is its internal structure: the single bright ring
visible at moderate resolution actually comprises an infinite sequence
of increasingly narrow sub-rings, each exponentially thinner than the
last.  Understanding this structure explains why the ray tracer
(\cref{sec:geo-raytracing}) must shrink its step size dramatically for
near-critical rays, and why the resulting images carry information about
the spacetime that no amount of astrophysical modelling can mimic.

\paragraph{Logarithmic winding.}

A photon approaching the Schwarzschild photon sphere at $r = 3M$ with
$b$ slightly above $b_c = 3\sqrt{3}\,M$
(\cref{eq:critical-impact}) winds through a deflection angle that grows
without bound as $b \to b_c$.  The rate of growth is controlled by the
Lyapunov exponent of the unstable circular orbit --- the exponential
rate at which neighbouring trajectories diverge from the bound orbit.

A radial perturbation $\delta r$ about the circular orbit at $r = 3M$
grows exponentially in affine parameter:
%
\begin{equation}
  \delta r(\lambda) \sim e^{\pm\gamma_\lambda\,\lambda} \,, \qquad
  \gamma_\lambda = \sqrt{-\tfrac{1}{2}\,V''_{\text{eff}}(3M)}
    = \frac{L}{9M^2} \,,
  \label{eq:photon-lyapunov}
\end{equation}
%
where $V''_{\text{eff}}(3M) = -2L^2/(81M^4)$ is the curvature of the
null effective potential (\cref{eq:veff-null}) at the photon sphere ---
the negative sign confirms the instability --- and $L$ is the angular
momentum of the photon.  Evaluating the angular velocity at $r = 3M$ in
the same affine parametrisation gives
$\d\varphi/\d\lambda\big|_{r=3M} = L/r^2 = L/(9M^2)$, exactly equal to
$\gamma_\lambda$.  Since the growth rate and the winding rate coincide,
perturbations grow by a factor of $e$ for every radian of azimuthal
winding --- i.e.\ the Lyapunov exponent per radian is unity.
\physnote{The Lyapunov exponent $\gamma_\lambda$ in affine parameter is
related to the coordinate-time Lyapunov exponent $\gamma_L$ of
\cref{sec:geo-raytracing} by a rescaling:
$\gamma_L = \gamma_\lambda\,(\d\lambda/\d t)$, evaluated on the
circular orbit.}

The total deflection angle for a ray with impact parameter $b$ close to
$b_c$ follows from inverting the exponential growth.  A photon that
approaches to within $\delta r$ of the photon sphere winds through
$\Delta\varphi \sim \ln(1/\delta r)$ radians; since $\delta r$ is
proportional to $b - b_c$ at leading order, the deflection diverges
logarithmically:
%
\begin{equation}
  \Delta\varphi \approx
    -\ln\!\left(\frac{b}{b_c} - 1\right) + \text{const} \,.
  \label{eq:deflection-divergence}
\end{equation}
%
The formula holds for $b \to b_c^+$ (rays that approach the photon
sphere from outside); as $b \to b_c^+$, the photon completes an
arbitrarily large number of orbits before escaping.
\physnote{The coincidence $\gamma_\lambda = \d\varphi/\d\lambda$ is
specific to the Schwarzschild photon sphere.  In the Kerr geometry, the
Lyapunov exponent and orbital frequency differ, and the growth per
radian becomes a function of the orbit parameters.}

\paragraph{The photon ring decomposition.}

A systematic classification of the ring structure follows from counting
how many times a photon crosses the equatorial plane during its close
approach to the photon sphere \cite{Johnson2020}.  The integer $n$
labels these crossings: $n = 0$ rays are the ``direct'' image with no
significant winding; $n = 1$ rays complete a half-orbit, creating the
``lensing ring''; rays with $n \geq 2$ form the ``photon ring'' proper.
Each sub-ring occupies a narrow annulus on the observer's sky centred on
the critical curve (the locus of $b = b_c$).

Since the deflection angle increases by $\pi$ for each additional
half-orbit, the impact parameters contributing to the $n$-th sub-ring
satisfy $\Delta\varphi \approx n\pi$, and
\cref{eq:deflection-divergence} gives
$b_n/b_c - 1 \sim e^{-n\pi}$.  The angular width $\delta b_n$ of the
$n$-th sub-ring therefore decreases exponentially with a universal
ratio:
%
\begin{equation}
  \frac{\delta b_{n+1}}{\delta b_n} = e^{-\gamma_{\text{ph}}} \,, \qquad
  \gamma_{\text{ph}} = \pi \quad \text{(Schwarzschild)} \,.
  \label{eq:demagnification}
\end{equation}
%
Each successive sub-ring is thinner by a factor of
$e^\pi \approx 23.1$.  The $n = 1$ ring --- the most prominent --- has
an angular width of a few microarcseconds for M87$^*$; the $n = 2$ ring
is $23$ times thinner, and $n = 3$ is over $500$ times thinner still.
The exponential narrowing is a direct consequence of the Lyapunov
instability: each additional half-orbit amplifies the radial deviation
by $e^\pi$, so rays completing one more half-orbit must have started
proportionally closer to the critical impact parameter.
\histnote{Darwin identified the photon sphere and logarithmic winding
in 1959 \cite{Darwin1959}.  The sub-ring decomposition and its
observational consequences were developed systematically by Johnson
et al.\ \cite{Johnson2020}.}

\begin{figure}[htbp]
  \centering
  % photon-ring-structure.tex — Photon ring sub-structure schematic
% Inputable TikZ fragment for fig:photon-ring-structure
% Requires: tikz with calc, arrows.meta, decorations.markings (loaded by ehbook.cls)
%
% Two-panel figure.
% Left: ray trajectories for n=0 (gold), n=1 (blue), n=2 (crimson)
%   around a Schwarzschild black hole (M=1).
% Right: observer's image plane showing sub-ring annuli.
%
\begin{tikzpicture}[
    >={Stealth[length=5pt]},
    font=\footnotesize,
  ]

  % ── Panel (a): Ray trajectories ──────────────────────────────────────

  \begin{scope}[shift={(0,0)}]

    % Scale: 0.55 cm per M.  Black hole r=2M, photon sphere r=3M.
    \def\scl{0.55}

    % Black hole (filled circle, r_+ = 2M)
    \fill[black] (0,0) circle ({2*\scl});

    % Photon sphere (dashed circle, r = 3M)
    \draw[EHMarginGray, dashed, thin] (0,0) circle ({3*\scl});
    \node[font=\scriptsize, text=EHMarginGray, anchor=south west]
      at ({3*\scl*cos(55)}, {3*\scl*sin(55)}) {$r\!=\!3M$};

    % Observer indicator (upper right)
    \draw[->, EHMarginGray, thin]
      (2.6, 1.8) -- ++(0.35, 0.25)
      node[right, font=\scriptsize] {observer};

    % ── n = 0 ray (gold): mild deflection ──────────────────────────────
    \draw[EHAccretionGold, thick, ->]
      plot[smooth, tension=0.55] coordinates {
        (-2.8, -1.0)
        (-1.2, -0.2)
        (0.2, 0.6)
        (1.5, 1.2)
        (2.6, 1.8)
      };
    \node[font=\scriptsize, text=EHAccretionGold]
      at (-2.8, -0.65) {$n = 0$};

    % ── n = 1 ray (blue): half-orbit winding ──────────────────────────
    \draw[EHJetBlue, thick, ->]
      plot[smooth, tension=0.5] coordinates {
        (-2.8, -0.15)
        (-1.6, 0.2)
        (-0.6, 0.9)
        (-0.2, 1.6)
        (0.3, 1.8)
        (1.1, 1.6)
        (1.8, 1.3)
        (2.6, 1.8)
      };
    \node[font=\scriptsize, text=EHJetBlue]
      at (-2.8, 0.2) {$n = 1$};

    % ── n = 2 ray (crimson): full orbit winding ───────────────────────
    % Trajectory hugs the photon sphere (r=3M, radius 1.65cm) more tightly
    \draw[EHHorizonCrimson, thick, ->]
      plot[smooth, tension=0.45] coordinates {
        (-2.8, 0.5)
        (-1.8, 0.8)
        (-1.1, 1.35)
        (-0.2, 1.7)
        (0.6, 1.65)
        (1.3, 1.15)
        (1.55, 0.4)
        (1.3, -0.35)
        (0.6, -0.95)
        (-0.2, -1.15)
        (-0.9, -0.85)
        (-1.2, -0.2)
        (-0.9, 0.5)
        (-0.2, 1.0)
        (0.7, 1.2)
        (1.6, 1.5)
        (2.6, 1.8)
      };
    \node[font=\scriptsize, text=EHHorizonCrimson]
      at (-2.8, 0.85) {$n = 2$};

    % Panel label
    \node[font=\scriptsize, anchor=north] at (0, -2.4) {(a) Ray trajectories};

  \end{scope}

  % ── Panel (b): Observer's image plane ────────────────────────────────

  \begin{scope}[shift={(5.8, 0)}]

    % Shadow: grey filled circle (radius = b_c on image plane)
    % Use 1.2 cm as the shadow angular radius on the image plane
    \def\bshadow{1.2}

    % Shadow fill
    \fill[EHMarginGray, opacity=0.12] (0,0) circle (\bshadow);
    \draw[EHMarginGray, thin] (0,0) circle (\bshadow);

    % n = 0 ring: wide annulus just outside shadow
    % Width ~ 0.55 cm (visually dominant direct image)
    \fill[EHAccretionGold, opacity=0.12]
      (0,0) circle ({\bshadow + 0.55}) -- (0,0) circle (\bshadow);
    \draw[EHAccretionGold, thick] (0,0) circle ({\bshadow + 0.55});
    % Re-draw inner boundary in gold
    \draw[EHAccretionGold, thick] (0,0) circle ({\bshadow + 0.01});

    % n = 1 ring: narrow annulus
    % Width ~ 0.55/23 ≈ 0.024 → exaggerate to 0.10 for visibility
    \def\rone{{\bshadow + 0.01}}
    \draw[EHJetBlue, thick] (0,0) circle ({\bshadow - 0.005});
    \fill[EHJetBlue, opacity=0.20]
      (0,0) circle (\bshadow) -- (0,0) circle ({\bshadow - 0.10});
    \draw[EHJetBlue, thick] (0,0) circle ({\bshadow - 0.10});

    % n = 2 ring: very narrow annulus (barely visible)
    \fill[EHHorizonCrimson, opacity=0.25]
      (0,0) circle ({\bshadow - 0.10}) -- (0,0) circle ({\bshadow - 0.14});
    \draw[EHHorizonCrimson, thick] (0,0) circle ({\bshadow - 0.14});

    % Labels for rings
    \draw[EHAccretionGold, thin, ->]
      ({(\bshadow + 0.28)*cos(35)}, {(\bshadow + 0.28)*sin(35)})
      -- ++(0.7, 0.45)
      node[right, font=\scriptsize, text=EHAccretionGold] {$n = 0$};

    \draw[EHJetBlue, thin, ->]
      ({(\bshadow - 0.05)*cos(145)}, {(\bshadow - 0.05)*sin(145)})
      -- ++(-0.6, 0.45)
      node[left, font=\scriptsize, text=EHJetBlue] {$n = 1$};

    \draw[EHHorizonCrimson, thin, ->]
      ({(\bshadow - 0.12)*cos(215)}, {(\bshadow - 0.12)*sin(215)})
      -- ++(-0.5, -0.5)
      node[left, font=\scriptsize, text=EHHorizonCrimson] {$n = 2$};

    % Shadow label
    \node[font=\scriptsize, text=EHMarginGray] at (0, -0.15)
      {shadow};
    \node[font=\scriptsize, text=EHMarginGray] at (0, -0.45)
      {$b = b_c$};

    % Demagnification annotation
    \draw[EHMarginGray, thin, <->]
      ({\bshadow + 0.55 + 0.08}, -0.9) -- ({\bshadow + 0.08}, -0.9);
    \node[font=\scriptsize, text=EHMarginGray, anchor=west]
      at ({\bshadow + 0.60}, -0.9) {$\delta b_0$};

    \draw[EHMarginGray, thin, <->]
      ({\bshadow + 0.08}, -1.25) -- ({\bshadow - 0.10 - 0.08}, -1.25);
    \node[font=\scriptsize, text=EHMarginGray, anchor=west]
      at ({\bshadow + 0.12}, -1.25) {$\delta b_1 \approx \delta b_0/23$};

    % Panel label
    \node[font=\scriptsize, anchor=north] at (0, -2.4) {(b) Image plane};

  \end{scope}

\end{tikzpicture}

  \caption[Photon ring sub-structure in the Schwarzschild geometry]{%
    Photon ring sub-structure in the Schwarzschild geometry.
    Left: ray trajectories classified by the number of equatorial
    crossings $n$ during the close approach to the photon sphere
    (dashed circle, $r = 3M$).  The $n = 0$ ray (gold) is deflected
    without significant winding; the $n = 1$ ray (blue) completes a
    half-orbit; the $n = 2$ ray (crimson) completes a full orbit.
    All three rays escape to the same distant observer.
    Right: the corresponding sub-rings on the observer's image plane,
    shown schematically with widths decreasing by a factor of
    $e^\pi \approx 23$ between successive rings.  The shadow boundary
    (grey fill) corresponds to $b = b_c = 3\sqrt{3}\,M$.}
  \label{fig:photon-ring-structure}
\end{figure}

\paragraph{The photon shell in Kerr.}

In the Kerr spacetime, the photon sphere generalises to a
\emph{photon shell} --- a three-dimensional region supporting
spherical photon orbits (orbits with constant Boyer--Lindquist radius)
at a range of radii.  The equatorial prograde and retrograde photon
orbit radii bound this range \cite{Bardeen1972}:
%
\begin{equation}
  r_{\text{ph}}^{\pm}
  = 2M\!\left(1 + \cos\!\left(
      \tfrac{2}{3}\,\arccos(\mp a_*)\right)\right) ,
  \label{eq:kerr-photon-radii}
\end{equation}
%
where $a_* = a/M$ is the dimensionless spin
(\cref{sec:kerr-isco}) and the upper sign gives prograde orbits.
At $a_* = 0$, both reduce to $r = 3M$.  At $a_* = 1$, the prograde
orbit drops to $r = M$ --- coinciding with the degenerate event
horizon --- while the retrograde orbit moves out to $r = 4M$.
Non-equatorial spherical photon orbits fill the interval
$r_{\text{ph}}^+ \leq r \leq r_{\text{ph}}^-$, parametrised
continuously by the Carter constant $\mathcal{Q}$
(\cref{eq:carter-constant}).

Each spherical photon orbit defines a pair of critical impact
parameters $(b_c(r), q_c(r))$, where
$q_c = \sqrt{\mathcal{Q}}/E$ is the reduced Carter parameter
(\cref{eq:carter-constant}).  As $r$ varies across the photon
shell, these trace out the \emph{critical curve} on the observer's
sky --- the boundary of the black hole shadow
\cite{Chandrasekhar1983}.  For Schwarzschild, the critical curve is a
circle of radius $b_c = 3\sqrt{3}\,M$; for Kerr, it is a
$D$-shaped contour whose flattened side faces the prograde direction,
reflecting the asymmetry between co-rotating and counter-rotating
photon orbits.
\xrefnote{The critical curve and shadow shape are rendered by the ray
tracer via the integration and termination pipeline of
\cref{sec:geo-raytracing,sec:geo-termination}.}

The exponential demagnification (\cref{eq:demagnification}) carries
over to Kerr, but the Lyapunov exponent $\gamma_{\text{ph}}$ becomes a
function of position around the critical curve.  Prograde photon orbits,
sitting closer to the horizon in a deeper gravitational potential,
have larger $\gamma_{\text{ph}}$ (faster demagnification); retrograde
orbits have smaller $\gamma_{\text{ph}}$.  For a rapidly spinning black
hole with $a_* = 0.9$, the demagnification ratio varies by roughly a
factor of three around the ring --- from about $50$ on the prograde side
to about $15$ on the retrograde side.
\implnote{The ray tracer's image implicitly contains all sub-ring
contributions: a photon that winds $n$ times is traced to its source
and rendered at the correct pixel.  No separate photon-ring mode is
needed --- the sub-ring structure emerges from the integration and
termination pipeline.}

\paragraph{Observational significance.}

The Event Horizon Telescope image of M87$^*$ resolves a bright ring of
diameter $\approx 42\;\mu\text{as}$, consistent with the $n = 0$ and
$n = 1$ contributions blurred together at the telescope's angular
resolution of $\sim 20\;\mu\text{as}$.  The $n = 1$ sub-ring carries a
disproportionate share of the ring brightness because it samples light
from the entire circumference of the black hole, acting as a
wide-angle gravitational lens.  Resolving the $n = 2$ ring --- which
would provide a direct, model-independent measurement of $\gamma_{\text{ph}}$ and
hence the black hole mass and spin --- requires baselines exceeding
Earth's diameter, motivating proposed space VLBI missions
\cite{Johnson2020}.
\physnote{In the eikonal (geometric optics) limit, the Lyapunov
exponent $\gamma_\lambda$ of the photon sphere determines the imaginary
part of the quasinormal mode frequencies \cite{Cardoso2009}.  The
photon ring structure is thus linked to the gravitational-wave
ringdown signal of a perturbed black hole.}

The photon ring decomposition reveals that the visually striking ring in
a black hole image is not one feature but infinitely many, nested with
each layer $e^\pi \approx 23$ times thinner than the last.  This
self-similar structure, set entirely by the Lyapunov exponent of the
photon sphere, is among the cleanest predictions general relativity
makes about observable black hole images.
